\chapter{Results and evaluation}
\section{Chapter overview}
Describe the evaluation types and how they validate your research outcomes.

\section{Quantitative results}
Present numerical metrics with clear explanation of trends and anomalies.

\begin{table}[H]
    \centering
    \begin{tabular}{|l|l|l|l|}
        \hline
        Model & Accuracy & F1 score & Dataset \\
        \hline
        Baseline & Insert value & Insert value & Insert dataset \\
        Proposed & Insert value & Insert value & Insert dataset \\
        \hline
    \end{tabular}
    \caption{Quantitative performance metrics (placeholder)}
    \label{tab:quantitative-results}
\end{table}

\section{Qualitative analysis}
Interpret visual outputs, case studies, and domain-expert observations.

\section{Ablation or sensitivity studies}
Analyze the impact of removing or changing major components.

\begin{table}[H]
    \centering
    \begin{tabular}{|l|l|l|}
        \hline
        Variant & F1 score & Remarks \\
        \hline
        Full model & Insert value & Reference setup \\
        Variant A & Insert value & Effect of change \\
        \hline
    \end{tabular}
    \caption{Ablation study results (placeholder)}
    \label{tab:ablation-results}
\end{table}

\section{Baseline or SOTA comparison}
Compare proposed approach with baselines and state-of-the-art methods using fair protocols.

\section{Discussion}
Interpret findings against research questions, objectives, and practical implications.

\section{Chapter summary}
Summarize key findings and transition to conclusion and future work.
